% !TEX root = ../main.tex
\chapter{System Model}
\label{ch:architecture}

% =====================================================================
% 骨架來源：docs/thesis/thesis-outline.md §3 + §4 + §6，併入舊 §10（泛化）的分層論述
%
% 【本章與第 4 章的分界】
%   第 3 章 = 靜態存在的東西：問題的形式化、規格長什麼樣、規格怎麼被執行
%   第 4 章 = 它怎麼被產生出來的：derive -> gate -> oracle -> repair
%
% 【結構調整說明】舊大綱把「泛化到 docx/xlsx」放成獨立的第 10 章最小 PoC。
% deck 已把主評測載體翻轉成 docx/xlsx，所以泛化不再是附加章節，
% 而是本章 §3.2 的分層論述 + 第 5 章的評測本體。DWG 降為最難的載體。
% =====================================================================

\section{Problem Formulation}
\label{sec:design:formulation}

\subsection{A Stateful, Two-Stage Input Model}
\label{sec:design:inputs}

The abbreviation $\mathcal{A}:T\rightarrow S$ is potentially misleading because the
system does not consume all experimental material in one call, and derivation is neither
pure nor deterministic. Let $d$ denote one derivation dataset in template lineage
$\ell$. We split the available material into a derivation layer and a later regression
layer,
\begin{align}
T &= \bigl(T_{\mathrm{der},d}^{(n)},T_{\mathrm{reg},\ell}\bigr), \\
T_{\mathrm{der},d}^{(n)}
  &= \bigl(d,\ell,B_d,Y_d^{?},\{X_{dj}\}_{j=1}^{N_d},R_d^{?},U_d^{?},
      Q_d,P,\Gamma,\Theta,H_d^{(<n)}\bigr), \\
T_{\mathrm{reg},\ell}
  &= \mathcal{R}_{\ell}
   = \left\{c_i=\bigl(I_i,\{X_{ij}\}_{j=1}^{N_i},u_i,\pi_i,Y_i^{?}\bigr)
     \right\}_{i=1}^{K_{\ell}} .
\end{align}
Here, $B_d$ is exactly one baseline template; $Y_d^{?}$ is an optional expected
output; $\{X_{dj}\}_{j=1}^{N_d}$ is a possibly empty collection of source artifacts;
$R_d^{?}$ is an optional operational-rule artifact; and $U_d^{?}$ is an optional
inline source-data override. When present, the override bypasses the locally parsed
artifact dictionary and, on the agent path, takes precedence over extracted keys. $Q_d$
groups the inspection level, force-refresh choice,
content-addressed dump, cache state, CAD preflight and health observations, locator
verification, and deterministic validation/repair feedback. The derivation environment
is also an input: $P$ denotes the prompt files, $\Gamma$ the agent's tool and skill
definitions, and $\Theta$ the model, gateway, and MCP configuration. Thus, changing a
prompt, an allowed tool, or the selected model changes $T_{\mathrm{der}}$ even when all
customer artifacts are byte-identical.

The term $H_d^{(<n)}$ is the pre-existing state before the $n$th derive. It contains
the field names and rule inventories accumulated from all existing non-degraded specs
derived from the same dataset. The implementation reads this state through
\texttt{derivation\_service.DerivationService.\_reference\_specs},
\texttt{\_field\_hints}, and \texttt{\_previous\_rule\_inventory}; the degraded flag
used by that selection is produced by \texttt{cad\_health.degradation}. The hints are
fed to the agent and the prior inventory is merged into the new spec evidence. With
$\omega_n$ denoting model sampling and the observed session/tool trace, the honest
derivation relation is therefore
\begin{equation}
S_d^{(n)} \sim
\mathcal{A}\!\left(T_{\mathrm{der},d}^{(n)};\omega_n\right),
\qquad
T_{\mathrm{reg},\ell}\notin\operatorname{dom}(\mathcal{A}).
\label{eq:stateful-derivation}
\end{equation}
Consequently, two derives over the same artifact bytes need not have the same context or
the same output: $H_d^{(<n)}$ may have changed, and $\omega_n$ is stochastic. In this
thesis, $\mathcal{A}:T\rightarrow S$ is only shorthand for the projection
$\mathcal{A}:T_{\mathrm{der}}\rightsquigarrow S$ in
Equation~\ref{eq:stateful-derivation}, not a claim that $\mathcal{A}$ is a pure function.

This split follows the implemented call boundary. The dataset creation and derive
paths, \texttt{api.create\_dataset} and
\texttt{derivation\_service.DerivationService.derive}, admit one baseline, at most one
expected artifact, $N_d$ source artifacts, and at most one rules artifact. They do not
read \texttt{models.DwgTestCase}. Only after $S$ exists does the separate endpoint call
\texttt{regression.RegressionRunner.run}, which selects every active test case of the
spec's lineage and executes them. $K_{\ell}$ is therefore the number of active cases at
regression time; the implementation imposes no $K_{\ell}=3$ limit. Each case $c_i$
contains an input template $I_i$, its own source artifacts and explicit source-data
override $u_i$, runtime parameters $\pi_i$, and an optional expected output $Y_i^{?}$.

The prompt, tools, and model terms above are concrete inputs rather than abstract
nuisance variables. \texttt{derivation\_agent.OpenCodeAgent.derive\_rules} selects the
\texttt{dwg-derivation} agent and \texttt{OpenCodeAgent.\_run} submits the task payload;
the selected agent is configured by \texttt{opencode.jsonc}, its prompt is
\texttt{prompts/derivation.md}, and its permitted CAD MCP tools are defined by the agent
configuration and \texttt{mcp\_server.server}. These files and settings must therefore
be versioned with the material definition used in an experiment.

We reason about the resulting specification as
\begin{equation}
S=\bigl(G,L_s,L_t,F,\Pi,\mathcal{E}_S\bigr),
\end{equation}
where $G$ contains value-source and transformation semantics, $L_s$ contains source
locators, $L_t$ contains independently derived target locators, $F$ contains rendering
rules, $\Pi$ is the runtime-parameter schema, and $\mathcal{E}_S$ is derivation and
validation evidence. The persisted \texttt{DwgSpec} stores these concerns across
\texttt{field\_map}, rule projections, \texttt{runtime\_params}, and
\texttt{evidence}; this tuple is a conceptual separation rather than a claim that each
component is a distinct database column.

\paragraph{Information boundary.}
The operational-rule artifact $R_d$ may identify source fields and state formulas,
conditions, constants, and formatting requirements, but it must not contain target cell
addresses, bookmark names, content-control tags, or mappings of the form
\texttt{source\_field $\rightarrow$ target\_cell}. Target references belong to $L_t$ and
must be derived from the template by $\mathcal{A}$; otherwise $R_d$ becomes an answer
table and the derivation problem is leaked into its input. \textbf{Status:
normative/aspirational only.} The boundary is stated but not enforced, in two
independent ways. First, no admission check exists: neither
\texttt{api.create\_dataset} nor
\texttt{derivation\_service.DerivationService.\_load\_rules} rejects an answer-bearing
rules artifact, so a hand-authored $R_d$ containing target addresses is accepted
silently. Second, the persisted projection is not structurally safe:
\texttt{derivation\_service.DerivationService.derive} copies the field's target-reference
slot into \texttt{DwgSpec.mapping\_rules}. That slot is carrier-agnostic --- its
interpretation follows the target kind of the template, so it holds bookmark identities
for \texttt{.docx} and cell addresses for \texttt{.xlsx}, which are exactly the two
artefact classes this boundary forbids. (The field retains the legacy name
\texttt{handles} from the drawing-oriented prototype; the naming is historical and does
not narrow the scope.) If that projection is exported, presented as the rule table, or
fed back as $R_d$, it exposes target answers, and $\mathcal{A}$ is then measured on its
ability to copy them rather than to derive them. It must therefore remain internal, and
be stripped of target references or split from the semantic rules, before the boundary
may be described as a property of the system rather than a requirement on it.

\subsection{Compile-Once, Execute-Many}
\label{sec:design:compile-once}

The target problem class consists of recurrent template-instantiation tasks for which
the output is deterministically derivable from structured source data and explicit
operational rules. ``Compile once'' means that the expensive, stochastic derivation in
Equation~\ref{eq:stateful-derivation} produces one reusable $S$; it does not mean that a
single fixed output job is reused, because instance values differ. For instance input
$V_i=(I_i,\{X_{ij}\},u_i,\pi_i)$, execution deterministically instantiates and applies
the spec,
\begin{align}
(J_i,P_i) &= \mathcal{C}(S,V_i), \\
D_i &= \mathcal{X}(J_i,I_i), \\
\operatorname{ok}_i
  &= \mathcal{O}_{\mathrm{plan}}(P_i,I_i,D_i)
     \land
     \begin{cases}
       \mathcal{O}_{\mathrm{expected}}(D_i,Y_i), & Y_i^{?}=Y_i,\\
       \mathrm{true}, & Y_i^{?}\text{ is absent}.
     \end{cases}
\end{align}
$\mathcal{C}$ is the deterministic per-instance compiler, $P_i$ is its authoritative
write plan, $\mathcal{X}$ is the format-specific executor, and $D_i$ is the produced
document. The problem is to derive an $S$ once such that $\operatorname{ok}_i$ holds
for every admissible instance, while the cost of derivation is amortized over repeated
executions.

The deterministic core is implementation-backed. In
\texttt{compiler.compile\_cad\_job}, the spec, source data, runtime parameters, and
current template values are explicit inputs; failures are returned before submission,
and the canonical job JSON is content-hashed. The DOCX/XLSX path shares this compiler
and applies $P_i$ through \texttt{docx\_target.DocxTargetAdapter.apply} or
\texttt{xlsx\_target.XlsxTargetAdapter.apply};
\texttt{document\_production.document\_job} records a
canonical hash for the document job. Regression is also implementation-backed:
\texttt{regression.RegressionRunner.\_run\_case} compiles each case, and
\texttt{regression.RegressionRunner.\_verify\_case} uses
\texttt{verification.verify\_output} plus
\texttt{verification.compare\_expected} when $Y_i$ exists.

This formulation deliberately localizes randomness to $\mathcal{A}$, and the execution
path now matches that localization. An earlier revision of this chapter recorded that it
did not: \texttt{production.ProductionController.start\_run} invoked
\texttt{derivation\_agent.OpenCodeAgent.repair\_cad\_job} after a failed gateway job
validation, and the comparison that guarded the repaired job excluded target locations,
so a model-relocated target could be accepted and submitted inside the same run. That
branch has since been removed from execution. A compiled job the tenant manifest rejects
terminates the run with \texttt{CAD\_JOB\_MANIFEST\_REJECTED} and an instruction to
re-derive, review, and republish; a cached target reference the current template no
longer contains terminates it with \texttt{CACHED\_HANDLE\_NOT\_FOUND}. Neither is
repaired in place, and \texttt{repair\_cad\_job} has no remaining caller on the execution
path.

\textbf{Status: implementation-backed, and pinned by a negative control rather than by
inspection.} \texttt{test\_successful\_production\_operation\_records\_zero\_llm\_calls}
replaces both agent repair entry points with a recorder that returns a failure, runs a
complete production operation, and then asserts both that the recorder was never invoked
and that the \texttt{LlmOperation} persisted for the run holds zero model calls. Reading
the source would only show that the call site is absent today; the test also fails if a
future change reintroduces one. This is the evidence behind claim C4, and it is
deliberately the narrowest of the thesis's claims: it covers $\mathcal{X}$, not
$\mathcal{A}$, and not the derivation-time repair loop of Chapter~\ref{ch:method}, which
is by construction inside the compile step whose cost repeated execution amortizes.

\subsection{Required Properties of the Specification IR}
\label{sec:design:properties}

The reusable IR must satisfy the following five properties. Each status below refers
to the repository state examined for this thesis, not merely to the intended design.

\paragraph{P1---separation of structure and instance values.}
$S$ must describe reusable source/target locators, transformations, formatting, and
parameter types; a particular instance supplies values through $V_i$. The only values
permitted in $S$ are constants explicitly prescribed by $R_d$, not values copied from a
calibration output. \textbf{Status: implementation-backed at the representation
interface.} \texttt{compiler.compile\_cad\_job} accepts \texttt{spec},
\texttt{source\_data}, and \texttt{runtime\_params} separately, while
\texttt{value\_layer.resolve\_field\_raw} evaluates the value layer at execution time;
\texttt{regression.RegressionRunner.\_run\_case} reuses the same spec with each case's
inputs.
This interface supports reuse, although successful regression---not the type shape
alone---is what detects a semantically overfit constant.

\paragraph{P2---explainable and verifiable locators.}
Every source or target locator must expose what it resolved to and must be checkable
against the corresponding artifact; an unresolved or ambiguous locator is evidence, not
permission to guess. \textbf{Status: implementation-backed.}
\texttt{source\_locator.resolve\_locator} returns the resolved sheet, cell, selection,
and match count, and \texttt{source\_resolution.verify\_against\_sample} checks that the
locator reproduces the derivation sample. On the target side,
\texttt{target\_ref.BaseTargetAdapter.verify} and \texttt{describe} provide a common
existence and explanation interface.
\texttt{docx\_target.DocxTargetAdapter.inspect} and
\texttt{xlsx\_target.XlsxTargetAdapter.inspect} retain labels, physical locations,
positional flags, and ambiguous references, while their \texttt{apply} methods return
every missed reference.

\paragraph{P3---bounded layout resilience and fail-loud behavior.}
Semantic identities should survive admissible layout changes, while a change outside
the declared tolerance must stop the run rather than trigger a guessed write. This
includes the zero-LLM execution requirement: an execution-time model must not relocate a
target. \textbf{Status: implementation-backed for the fail-loud half; the resilience
half is bounded and measured rather than general.} Source locators resolve by
sheet/header/section semantics; \texttt{docx\_target.DocxTargetAdapter} distinguishes
stable content-control/bookmark identities from positional table cells;
\texttt{xlsx\_target.XlsxTargetAdapter} distinguishes stable defined names from
positional coordinates; and both reject duplicate or missing references. The execution
branch that previously undermined this property --- agent repair of a rejected job
inside the run --- has been removed, as described in
Section~\ref{sec:design:compile-once}, so a target that cannot be resolved now ends the
run with a code instead of being relocated.

What is \emph{not} claimed is the resilience half in general form. Layout tolerance is
declared per locator mechanism, and a change outside a mechanism's declared tolerance is
detected only insofar as the mechanism can see it. Two measured boundaries are carried
forward to Section~\ref{sec:design:boundaries} rather than glossed here: an unused
detail-row cell that no column of a repeat block addresses is not cleared by
\texttt{unused\_slots: "blank"}, so a template reused from a previous reporting period
can carry a stale value into a fresh run; and the derivation-time semantic oracle
compares only references whose text differs between the baseline and the expected
document, so a specification that writes into a cell the expected document leaves blank
is structurally invisible to it.
\paragraph{P4---diffability, versioning, and lineage.}
The IR and its evidence must support reviewable changes, immutable publication, and
regression over a template lineage. \textbf{Status: implementation-backed for the
currently produced handle-backed field-map shape.}
\texttt{models.DwgSpec} carries a lineage-scoped unique version and template fingerprint;
\texttt{drift.spec\_field\_drift} reports added, removed, and attribute-level field
changes; \texttt{models.DwgProductionVersion.save} prevents mutation of the published
spec/hash identity; and \texttt{regression.RegressionRunner.run} persists results for
all active cases of the lineage. The drift function reads \texttt{target.handles} rather
than the generic \texttt{targets} representation; diffability for a future spec that
uses only generic target refs remains normative until that path is implemented.

\paragraph{P5---independently auditable rule coverage.}
$S$ must enumerate which rules and fields it claims to cover, and the denominator must
come from the rule artifact independently of the model being scored. \textbf{Status:
implementation-backed since the coverage-v2 path replaced agent self-enumeration.} The
original implementation --- \texttt{rule\_coverage.merge\_inventory} over an agent-supplied
\texttt{rule\_inventory}, unioned across prior derives --- did not satisfy P5 and is
retained only to read historical specifications. A rule the agent never enumerated was
absent from the denominator, so ``22 of 22 covered'' and ``22 of 40 covered'' were the
same observation.

The active path externalizes the denominator instead. A dataset that carries a rule
document must also carry a frozen, versioned rule-applicability manifest;
\texttt{rule\_coverage.parse\_rule\_manifest} verifies that the manifest names the exact
SHA-256 of the rule-document bytes and that its declared document roles match the
dataset's, and \texttt{derivation\_service.DerivationService.\_load\_rule\_manifest}
refuses the derive with \texttt{RULE\_MANIFEST\_MISSING} before any paid model call, so a
run cannot discover after the fact that its release denominator is unauditable. The agent
supplies a \emph{disposition} for each externally assigned rule identifier and nothing
else: \texttt{rule\_coverage.score\_manifest\_coverage} discards unknown identifiers,
refuses duplicate dispositions, and refuses an assignment whose document roles disagree
with the manifest, so the scored agent can neither add, remove, rename, nor move a rule
between roles. The resulting partition into \texttt{covered}, \texttt{unresolved},
\texttt{rejected}, and \texttt{unmentioned} is mutually exclusive and sums to the frozen
applicable set, and \texttt{gates.canonical\_unresolved} turns every non-\texttt{covered}
disposition into a publish gate rather than partial credit. Section~\ref{sec:der:coverage}
develops what this measures and what it still does not.
\subsection{Measured Boundaries of the Specification IR}
\label{sec:design:boundaries}

The five properties above state what the IR must do. This subsection states what it
demonstrably does not, because a boundary that is only discovered by a reader running the
system is indistinguishable from a defect the author did not notice. Each item below is
either pinned by a test or was measured on a real artifact; where an alternative shape was
considered and rejected, the reason is given, since ``we did not think of it'' and ``we
rejected it'' are different claims.

\paragraph{B1---\texttt{exclude} is a row-content predicate.}
A range's \texttt{exclude} clause inspects the content of a data row. It does not inherit
meaning from a preceding section header, so a sheet whose rows alternate between
``current month'' and ``cumulative'' groups under one header cannot be partitioned by
\texttt{exclude}. The rejected alternative was a domain flag such as
\texttt{exclude\_cumulative}; it was refused because it names one agency's form inside a
carrier-neutral IR. The general shape adopted instead is a group projection
(\texttt{records}: a fixed \texttt{group\_size} with named row roles and per-column
\texttt{select\_offset}), which describes the same layout without naming the domain.

\paragraph{B2---stacked records assume a fixed group size.}
\texttt{records} models ``$g$ physical rows form one logical record'' for a constant $g$.
A record whose row count varies with its own data is outside the IR. Both the read side
and the write side are implemented, and a specification produced by the agent for the
XLSX evaluation template declares \texttt{records} with \texttt{group\_size: 2} and named
row roles, so the primitive is exercised by a main evaluation dataset rather than by
fixtures alone.

\paragraph{B3---the verification basis for the collection primitives is $n=1$.}
\texttt{artifact\_collection}, \texttt{records}, and the equality merge of
Section~\ref{sec:design:loc:range} each have exactly one motivating real workflow. Six
independently sourced government workbooks were scanned for a second instance of a
repeating multi-row record block and none was found; the nearest candidate was a one-off
label merge, not a repeating record. The generalization threshold this thesis sets for a
new primitive is therefore not met for these three, and they are presented as
carrier-neutral shapes with a single validating workflow rather than as generally
validated primitives. The risk this creates is specific and worth naming: it does not
appear as a failing test, because the tests were written from the same single example.
It appears as a wrong result on a form nobody has tried.

\paragraph{B4---merge is equality-only and inner.}
Two resolved row sets can be joined on equal key columns. Range containment, interval
overlap, and inequality correspondences are not expressible. The join is deliberately
inner: supplying a default for an unmatched key would convert a missing record into a
normal-looking number without anyone having authorized the default, which is precisely the
silent-failure shape this thesis measures. A specification that needs the filled row must
produce it on the source side, where the decision is visible. Merges chain, so three
tables can be joined as two binary merges; the agent produced a chained pair unprompted.

\paragraph{B5---the aggregation axis is columns only.}
Aggregation sums named columns. The criterion is who decides the number of columns: if the
template does, naming them is legitimate; if \emph{this instance's data} does, the
specification has become instance-specific, and the correct modelling is a long table plus
a pivot at write time. A row axis is not offered because a wide table's column names would
then be data values, which would disable the author-time static check that a formula only
names columns the row set actually has.

\paragraph{B6---a repeat block clears only the cells its own columns address.}
When a block declares \texttt{unused\_slots: "blank"}, the cells cleared in an unused slot
are the cells that block's columns point at. A cell in the same slot that no column
addresses is not cleared, because the IR knows a block's targets only through the strings
its columns declare, and ``every other cell on this row'' is not something it can name.
For a pristine template this is exactly right, and every instance of both evaluation
datasets starts from a pristine template. For the common practice of reusing last period's
completed report as this period's template, it means a stale value can sit on a row the
system considers written. The behavior is pinned by a test rather than left to inference,
and the rejected workaround---declaring a column whose value is always blank---was refused
because it puts layout knowledge into the value layer and makes the column count grow with
the template's width.

\paragraph{B7---one instance is one filled target template.}
A specification may declare several output documents, but their roles must be distinct; a
run that would produce $M$ documents of the same role is refused. Relaxing this is a
separate decision and is not claimed here.

\subsection{Design Principle: Fail Loud, Not Silent}
\label{sec:design:failloud}
% TODO ★ 尚未動筆 ★ 這一小節在 2026-08-15 撰寫 §3.1 時被整段刪除，同日復原為寫作規格。
%      刪除它不是決定，是疏漏——本節是全篇的核心設計信條，不可省略。
% TODO 明確寫成本論文的核心設計信條，並回指第 2 章的失敗模式分類。
%      要點：一個有理有據但錯的輸出，比一個明白失敗的輸出更危險，因為它會被交付出去。
%      §3.1.1 的 information boundary 與 §3.1.3 的 P3、P5 都是這條信條的實例，
%      寫的時候要把它們收攏成同一個原則，不要各講各的。

\subsection{Scope and Applicability Conditions}
\label{sec:design:scope}
% TODO ★ 收窄後的命題寫在這裡，Introduction 只放結論句 ★
%      2026-08-15 註：命題本身已寫進 §3.1.2（Compile-Once, Execute-Many）的正文，
%      但那裡是順著 compile-once 的論證帶出來的，沒有獨立成「適用條件」的判準清單，
%      也沒有往第 6 章的前向引用。本小節仍要寫，並與 §3.1.2 交叉引用而非重複。
%      適用條件（三條都要成立）：
%        1. recurrent —— 同一份模板會被實例化多次（否則過不了 break-even）
%        2. deterministically derivable —— 輸出可由結構化來源資料確定性推導
%        3. explicit operational rules —— 存在明確的操作規則文件，而非需要真實推理的判斷
%      不適用的情境在第 6 章 §6.2 展開，這裡只給定義不談後果。

% =====================================================================

\section{Domain-Dependent and Domain-Independent Layers}
\label{sec:design:layers}

% TODO 這一節決定論文被當成 tool paper 還是 method paper，要放在系統描述的最前面。
%      領域相關：dump 工具、locator 的實作、target adapter
%      領域無關：IR 結構、gates、oracle 種類、repair loop、lineage、覆蓋率計分
%      三條載體（docx / xlsx / dwg）共用同一條 derive -> gate -> compile -> run 管線。

\section{Specification IR}
\label{sec:design:ir}

% TODO 欄位定義、locator、值來源、格式化規則、約束。

\section{Structured Dump as Grounding}
\label{sec:design:dump}

% TODO docx / xlsx 的 inspect 與傾印；CAD 端的 DumpDwg 與 labeled-field 解析。
% TODO 【誠實揭露】openpyxl 的 load/save 有量測到的損失（儲存格註解、列印設定），
%      python-docx 亦有邊界情形。這些屬於載體限制而非方法限制，但要寫出來。

\section{Locator Mechanisms}
\label{sec:design:locators}

\subsection{Semantic Locators}
\label{sec:design:loc:semantic}

% TODO label -> value 的鄰近／對齊關係；以 section 限定搜尋範圍避免跨區塊誤命中；
%      sheet alias 與 fallback 策略。

\subsection{Range Locators and Source-Side Aggregation}
\label{sec:design:loc:range}

% TODO 【寫作規格 2026-08-23 更新——原註記說「實作尚未開始」，那是錯的，
%      而且方向是低估系統。ADR 0002 是 2026-08 以後每一份規格的地基，
%      它的讀取端、顯式邊界、來源聚合與 evidence 都在 production path 上。
%      照舊註記寫會把已完成的能力寫成未來工作。】
%
%      本小節要寫的是**一條演進線**，不是一個機制，因為每一步都是被實測逼出來的：
%        1. ADR 0002（Accepted，G1–G4／G6 已實作；G5 是 benchmark case assertion）
%           範圍定位器 ＋ 來源端聚合。約束：外延必須顯式宣告，不允許隱式整欄。
%           難題是外延：明細第幾列到第幾列、含不含小計列、插列後還對不對。
%           外延差一列會產生「有理有據但錯」的數字，是沉默失敗的教科書範例。
%        2. ADR 0004／0005：把一段列壓成純量不夠——spec 只觸及模板 14 個明細列位的
%           前 3 個，因為校準實例只有 3 筆。那不是 compile-once 是 compile-per-instance。
%           修法是列投影 ＋ repeat 區塊 ＋ 具名外延 ＋ 多文件。
%        3. ADR 0006：repeat 的欄只能直接讀一個來源欄，做不到逐列計算與跨檔查表。
%           修法是**外延先解析成 row set**，逐列公式、join、合計都在那張表上做。
%        4. ADR 0009：`{slot}` 既沒有起點也沒有間距，所以資料區不從第 1 列開始的
%           工作表定不到位。修法是 slot_origin + records（group_size / row_role）。
%        5. ADR 0010：兩個 row set 的等值合併。
%      這條線本身就是論文的論述：**表達力邊界是被資料集逼著往外推的，不是設計出來的**。
%
%      表達力邊界（寫進 §3.9，不要在這裡展開）：等值合併不做 outer join、不做範圍型
%      對應；聚合軸只有欄向；records 假設 group_size 固定；ADR 0008／0009／0010 的
%      驗證基礎都是 n=1。

\subsection{Failure Modes and What Each Mechanism Blocks}
\label{sec:design:loc:coverage}

% TODO 做成一張表，橫軸是第 2 章的失敗模式，縱軸是各機制。
%      這張表證明設計不是拼湊出來的。
% TODO 【現況要誠實寫】最近一次 derive 對映 31 條規則，21 條無法確認寫入位置，
%      成因是語意的（target_not_found / target_ambiguous）。這個數字要嘛先修掉，
%      要嘛作為 locator 機制現階段限制寫進 §3.9。不可以靜靜略過。

\section{Target Reference and Write-Back}
\label{sec:design:targets}

% TODO docx（bookmark / content control）、xlsx（具名範圍 / 座標）、dwg（handle）
%      三種目標參照的統一契約，以及各自的可寫入性判定。

\section{Deterministic Compiler and Zero-LLM Runtime}
\label{sec:design:runtime}

% TODO spec -> 可執行計畫；錯誤在編譯期而非執行期爆。受限 AST 求值。
% TODO value layer：實例輸入的驗證與型別。
%
% TODO ★ 論文級紅線，寫死在這一節 ★
%      執行期不得有 LLM fallback。related-work v1 審查判定 C4（zero-LLM execution）
%      是唯一乾淨的 claim，所以這不是實作偏好而是命題的必要條件。
%      具體推論：locator 定位失敗時必須失敗並回報，
%      不可以「叫 LLM 猜一個」來降低 §3.5.3 那個失敗率。

\section{Lifecycle Management}
\label{sec:design:lifecycle}

\subsection{Lineage and Versioning}
\label{sec:design:lineage}

% TODO spec 版本、產出版本、輸入版本三者的可追溯關係。

\subsection{Regression Suite}
\label{sec:design:regression}

% TODO 每次 spec 變更自動重跑全部案例。

\subsection{Drift Detection}
\label{sec:design:drift}

% TODO 模板改版時如何被發現、如何定位、如何最小化重新推導範圍。
%      fingerprint 對 docx dump 已驗證可用（文字改 -> 指紋不變、bookmark 改名 -> 指紋變）。
% TODO 【已知缺陷，修掉或寫出來】drift 目前只讀 target.handles 不讀通用的 targets，
%      一旦有 spec 使用通用 refs，漂移偵測會直接瞎掉。

\subsection{Auditability}
\label{sec:design:audit}

% TODO 為什麼這個架構能被工程流程接受，而 agent-per-instance 不能。
%      措辭注意：用 traceability / reproducibility / change control，
%      不要宣稱「法規要求 deterministic output」（找不到學術來源）。

\section{Expressiveness Boundary}
\label{sec:design:boundary}

% TODO 老實寫現階段表達不了的東西，這是第 6 章 future work 的來源：
%      - 寫入端不定長度明細（來源 7 筆就插 7 列）—— 插列會讓所有 positional ref 位移，
%        連動目標轉接器、比對、清空、fingerprint、drift。
%        替代形狀是「模板固定上限明細列位，未用列必須清空」，
%        而這正好把 negative obligation 變成一等公民
%      - 需要跨文件推理而非查表的欄位
%      - 【§3.5.3 的 locator 失敗率若沒修掉，也要寫進這裡】
